\section{Implications for Practice}

The clearest practical implication is that adopting a more structured, multi-agent workflow is not primarily a tooling decision. Google Antigravity supported sequential, role-specialized agents, but the actual source of Part 2's improvement was the Charter and the review structure built on top of that support --- a set of project-specific rules and gates that had to be authored, and that did the substantive work regardless of which environment executed them. An organization considering a similar move should expect to invest in writing and maintaining this kind of rule set as a first-class deliverable, not treat it as a byproduct of switching to a more capable agent platform.

A second implication follows from where the Charter's rules came from. Not all of it needs to be learned from scratch: a rule set can reasonably start from architectural conventions that are already widely accepted --- established patterns like a single application factory, thin route handlers, or a consistent service-layer error format are not specific to this project and could be adopted upfront from general software-engineering practice or from rules carried over from earlier projects. What this project adds on top of that starting point is the second source the Charter actually drew on: rules written directly from mistakes observed during Part 1. For a new project, this suggests treating a rule set as built in two layers --- established conventions adopted from the outset, and project-specific rules added as concrete failure modes are discovered during development --- rather than either writing everything in advance or waiting to learn everything the hard way.

A third implication concerns where oversight is placed rather than how much of it exists. The comparison in this thesis indicates that moving a human's approval earlier --- to a plan, before implementation, rather than to a finished result --- was more consequential than any change in how much the human reviewed. For practice, this reframes human-in-the-loop design as a question of sequencing: deciding at which stage of an agent's workflow a human is required to act, rather than only deciding whether review happens at all.

